%%%%%%%%%%%%%%%%%%%%%%%%%%%%%%%%%
% Methodology                   %
%%%%%%%%%%%%%%%%%%%%%%%%%%%%%%%%%

\section{METHODOLOGY}
\label{sec:methodology}

\noindent
This section outlines the systematic approach adopted in this study, detailing the research framework, the data generation and processing pipeline, and the tools and technologies employed to achieve the stated objectives.

\subsection{Design Science Research Methodology}

\noindent
Design Science Research (DSR) is a research paradigm whose primary aim is to develop and evaluate purposeful artifacts that address clearly defined practical problems \cite{wieringa2014design}. Unlike observational or explanatory research, DSR is prescriptive: it produces artifacts — models, methods, systems, or frameworks — and evaluates them against the problem context for which they were designed. This makes DSR particularly well-suited to applied computer science and engineering research, where the goal is not only to understand a phenomenon but to build something that improves upon the current state.

\par \vspace{0.3cm}
In this work, the artifact under development is a deep learning framework for synthesizer parameter estimation: a system that takes an audio signal and a MIDI pitch value as inputs and produces a vector of predicted synthesizer parameters as output. The problem context is the inverse synthesis problem — recovering the parameters of a software synthesizer from its audio output — which is a practical challenge faced by sound designers and music producers. DSR is applied here because the work involves designing and iteratively refining a computational artifact, evaluating it against quantitative performance criteria, and drawing design knowledge from the evaluation results to guide subsequent iterations.

\par \vspace{0.3cm}
The iterative nature of DSR aligns naturally with the three-phase experimental structure of this project. Each iteration — V1, V2, and V3 — corresponds to a complete DSR cycle in which a version of the artifact is designed, implemented, and evaluated, with findings feeding directly into the design of the next iteration. This structure ensures that each design decision is motivated by empirical evidence rather than assumption.

\subsection{Design Science Research Iterations}

The DSR process is structured around five interrelated phases, as illustrated in Figure \ref{fig:dsr_cycle}. These phases are not strictly sequential; rather, they form a cycle in which evaluation results inform subsequent problem investigation and solution redesign. The following subsections describe how each phase maps to the work conducted in this study.

\begin{figure}[H]
\centering
\begin{tikzpicture}[font=\small, >=Stealth]

\draw[line width=14pt, draw=gray!25] (0,0) circle (3.0cm);
\draw[-{Stealth[length=6mm]}, line width=14pt, draw=gray!45]
      (155:3.0cm) arc[start angle=155, end angle=110, radius=3.0cm];

\tikzset{
  dsrbox/.style={rounded corners=8pt, minimum width=4.1cm,
                 minimum height=1.2cm, align=center, font=\normalsize}
}

\node[dsrbox, fill=orange!80, text=white]   (prob)   at (90:3.0cm)   {Problem\\Investigation};
\node[dsrbox, fill=teal!65, text=white]     (design) at (20:3.0cm)   {Solution Design};
\node[dsrbox, fill=teal!50, text=white]     (valid)  at (-35:3.0cm)  {Design Validation};
\node[dsrbox, fill=blue!55, text=white]     (impl)   at (-110:3.0cm) {Implementation};
\node[dsrbox, fill=blue!70, text=white]     (eval)   at (180:3.0cm)  {Evaluation};

\end{tikzpicture}
\vspace{0.4cm}
\caption{Design Science Research cycle with the five phases applied in this study.}
\label{fig:dsr_cycle}
\end{figure}

\noindent Table~\ref{tab:dsr_mapping} maps each DSR phase to the specific activities carried out in this study, the artifact or output produced, and the section in which it is reported.

\begin{table}[H]
\centering
\renewcommand{\arraystretch}{1.4}
\begin{tabular}{p{2.8cm} p{4.8cm} p{3.6cm} p{2.0cm}}
\hline
\textbf{DSR Phase} & \textbf{Activity in this study} & \textbf{Output} & \textbf{Reported in} \\
\hline
Problem Investigation  & Literature review; gap analysis; formulation of research questions RQ1--RQ3 & Research questions; literature gap table & Sections 1, 2 \\
Solution Design        & Dataset strategy; input representation; architecture and training objective design & V1, V2, V3 design specifications & Section 4 \\
Design Validation      & Architecture correctness checks; V2 controlled ablation study isolating each component & Per-component MAE impact; ablation results & Sections 4, 5 \\
Solution Implementation & V1$\to$V2$\to$V3 model training across cross-validation splits & Trained model checkpoints & Section 4 \\
Solution Evaluation    & MAE and MSS evaluation on held-out test sets; per-parameter breakdown; architecture comparison & Results tables; analysis & Section 5 \\
\hline
\end{tabular}
\vspace{0.3cm}
\caption{Mapping of the five DSR phases to project activities, outputs, and report sections.}
\label{tab:dsr_mapping}
\end{table}

\subsubsection{Problem Investigation}

The problem investigation phase begins with a thorough analysis of the existing literature on inverse synthesis, audio deep learning, and synthesizer parameter estimation. This review, presented in Section 2, revealed several gaps in the current state of the art: (i) no prior work has applied deep learning-based parameter estimation to a modern wavetable synthesizer such as Vital; (ii) existing approaches rely on a single mel-spectrogram scale, ignoring the complementary information available at different time-frequency resolutions; (iii) no prior work incorporates explicit pitch conditioning to disentangle timbral features from the fundamental frequency; and (iv) the Audio Spectrogram Transformer and ConvNeXt architectures have never been directly compared in the context of synthesizer sound matching.

\par \vspace{0.3cm}
These gaps, combined with the practical motivation of building a tool useful to sound designers and producers, constitute the ``wicked problem'' that drives this research: a high-dimensional, non-linear, ill-posed inverse mapping whose difficulty is compounded by parameter interdependence, timbral ambiguity, and the non-differentiable nature of the target synthesizer. The three research questions formulated in Section 1 directly address the identified gaps and guide the design effort.

\subsubsection{Solution Design}

The solution design phase translates the problem formulation into a concrete artifact specification. Three core design decisions structure this phase.

\par \vspace{0.3cm}
First, the \textbf{dataset design}: rather than sampling parameters uniformly at random — which would produce mostly silent or musically incoherent sounds — a timbre-constrained generation strategy is devised. Parameters are sampled within musically meaningful sub-ranges defined for five timbre categories (bass, lead, pad, percussion, and random), and audio is rendered programmatically via the Pedalboard library driving Vital. The parameter space is restricted to 16 controls inspired by the canonical analog signal chain of the Minimoog and Prophet-5.

\par \vspace{0.3cm}
Second, the \textbf{input representation design}: a multi-scale mel spectrogram is proposed as a three-channel input tensor, combining fine, mid, and coarse time-frequency resolutions to provide the model with complementary views of each audio clip. This design directly addresses the resolution trade-off inherent in single-scale representations.

\par \vspace{0.3cm}
Third, the \textbf{architecture design}: three backbone families are selected for comparison — a custom residual CNN, a pre-trained ConvNeXt, and a pre-trained Audio Spectrogram Transformer — each augmented with a FiLM-based pitch conditioning branch and a grouped prediction head that assigns a dedicated sub-MLP to each semantically coherent parameter group.

\subsubsection{Design Validation}

Design validation is carried out through a systematic ablation study in the V2 iteration. Each architectural innovation introduced in V2 — multi-scale input, FiLM pitch conditioning, and the grouped prediction head — is evaluated in isolation by training variants in which only one component is removed or simplified at a time, while all other factors are held constant. This controlled approach ensures that performance differences are attributable to the specific component under test, rather than to confounding factors. The results of this ablation study, reported in Section 5, provide the empirical basis for the design choices carried forward into V3.

\subsubsection{Solution Implementation}

Implementation proceeds across three successive iterations, each building on the validated design of its predecessor.

\par \vspace{0.3cm}
\textbf{V1} establishes the baseline: a dataset of single-scale mel spectrograms is generated, and the three model families are trained and compared under identical conditions. The optional sinusoidal pitch embedding branch is evaluated as a simple concatenation-based conditioning mechanism.

\par \vspace{0.3cm}
\textbf{V2} implements all three architectural innovations simultaneously: multi-scale mel input, FiLM pitch conditioning, and the grouped prediction head with weighted loss. A full ablation suite is trained alongside the main models to validate each component.

\par \vspace{0.3cm}
\textbf{V3} addresses the key failure mode identified in V2 — the ineffectiveness of pitch conditioning when each preset is rendered at only a single MIDI pitch. By re-rendering every preset at three distinct MIDI notes per timbre, the dataset grows substantially, and the model is now exposed to the same parameter vector at different spectral positions, providing a meaningful conditioning signal for the FiLM branch. SpecAugment regularization is also introduced at this stage.

\subsubsection{Solution Evaluation}

Evaluation is conducted at two levels. The primary metric is the Mean Absolute Error (MAE) computed in the normalized $[0,1]$ parameter space, reported as mean $\pm$ standard deviation across five cross-validation splits to ensure statistical robustness. Per-parameter MAE is additionally reported to identify which synthesis controls are well-predicted and which remain structurally ambiguous. For V2 and V3, a Multi-Scale Spectral (MSS) loss is computed on a held-out test subset by re-rendering the predicted parameters through Vital and measuring the spectral distance between the original and reconstructed audio, providing a perceptual evaluation complementary to the parameter-space MAE.

% ─────────────────────────────────────────────────────────
\subsection{Data Generation Pipeline}

The data generation pipeline transforms raw synthesizer parameter configurations into a normalized, split dataset ready for model training. It consists of six sequential steps, illustrated in Figure~\ref{fig:pipeline}. Each step is designed to be checkpoint-aware and resumable, as the full pipeline involves several hours of audio rendering and processing.

\begin{figure}[H]
\centering
\begin{tikzpicture}[
  node distance=0.55cm and 0.9cm,
  box/.style={rectangle, rounded corners=6pt, draw=gray!60, fill=blue!7,
              text width=3.0cm, align=center, minimum height=1.3cm, font=\small},
  arrow/.style={-{Stealth[length=3.5mm]}, thick, gray!60}
]

% Row 1
\node[box] (s1) {\textbf{Step 1}\\[3pt]Parameter\\Extraction};
\node[box, right=of s1] (s2) {\textbf{Step 2}\\[3pt]Preset\\Generation};
\node[box, right=of s2] (s3) {\textbf{Step 3}\\[3pt]Audio\\Rendering};

% Row 2 (right to left below row 1)
\node[box, below=of s3] (s4) {\textbf{Step 4}\\[3pt]Silence\\Filtering};
\node[box, left=of s4]  (s5) {\textbf{Step 5}\\[3pt]Mel Spectrogram\\Extraction};
\node[box, left=of s5]  (s6) {\textbf{Step 6}\\[3pt]Dataset Split\\\& Normalization};

% Arrows
\draw[arrow] (s1) -- (s2);
\draw[arrow] (s2) -- (s3);
\draw[arrow] (s3) -- (s4);
\draw[arrow] (s4) -- (s5);
\draw[arrow] (s5) -- (s6);

% Corner connector
\draw[arrow] (s3.south) -- (s4.north);

\end{tikzpicture}
\vspace{0.3cm}
\caption{The six-step data generation pipeline, from parameter extraction to normalized dataset splits.}
\label{fig:pipeline}
\end{figure}

% ─────────────────────────────────────────────────────────
\subsubsection{Step 1: Parameter Extraction from Vital}

All 16 target parameters are enumerated programmatically from Vital via the Pedalboard Python library, which exposes the synthesizer's internal parameter table. For each parameter, the extraction records its name, internal plugin index, valid range $[x^{\min}_k,\, x^{\max}_k]$, and discretization information where applicable. Seven parameters — the four envelope times, the filter cutoff, the unison voice count, and the oscillator 2 transpose — are internally quantized by Vital to a finite set of valid discrete values. During preset generation, any continuously sampled value for these parameters is snapped to the nearest valid discrete point. The extracted ranges and discrete value tables are stored as JSON reference files and used throughout the pipeline to ensure that all generated presets are valid within Vital's internal representation. Table~\ref{tab:parameters} provides the complete listing of all 16 parameters in their prediction order, together with their group assignment, the normalization applied to raw Vital values before min-max scaling to $[0,1]$, whether the parameter is internally discretized by Vital, and the loss weight assigned to each group during training (see Equation~\ref{eq:weighted_loss} in Section 4).

\begin{table}[H]
\centering
\renewcommand{\arraystretch}{1.35}
\begin{tabular}{clp{2.2cm}p{2.4cm}cc}
\hline
\textbf{\#} & \textbf{Parameter} & \textbf{Group} & \textbf{Normalisation} & \textbf{Discrete} & \textbf{Loss wt.} \\
\hline
0  & \texttt{osc1\_wave\_frame}    & Oscillator & Linear        & ---  & $2\times$ \\
1  & \texttt{osc1\_level}          & Oscillator & Linear        & ---  & $2\times$ \\
2  & \texttt{osc1\_unison\_voices} & Oscillator & Linear        & Yes  & $2\times$ \\
3  & \texttt{osc1\_unison\_detune} & Oscillator & Linear        & ---  & $2\times$ \\
4  & \texttt{osc2\_wave\_frame}    & Oscillator & Linear        & ---  & $2\times$ \\
5  & \texttt{osc2\_level}          & Oscillator & Linear        & ---  & $2\times$ \\
6  & \texttt{osc2\_transpose}      & Oscillator & Linear        & Yes  & $2\times$ \\
\hline
7  & \texttt{env1\_attack}         & Envelope   & $\log_{10}$   & Yes  & $1\times$ \\
8  & \texttt{env1\_decay}          & Envelope   & $\log_{10}$   & Yes  & $1\times$ \\
9  & \texttt{env2\_attack}         & Envelope   & $\log_{10}$   & Yes  & $1\times$ \\
10 & \texttt{env2\_decay}          & Envelope   & $\log_{10}$   & Yes  & $1\times$ \\
\hline
11 & \texttt{filter1\_cutoff}      & Filter     & $\log_{10}$   & Yes  & $1\times$ \\
12 & \texttt{filter1\_resonance}   & Filter     & Linear        & ---  & $1\times$ \\
\hline
13 & \texttt{sample\_level}        & Misc       & Linear        & ---  & $1\times$ \\
14 & \texttt{mod1\_amount}         & Misc       & Linear        & ---  & $1\times$ \\
15 & \texttt{reverb\_mix}          & Misc       & Linear        & ---  & $1\times$ \\
\hline
\end{tabular}
\vspace{0.3cm}
\caption{All 16 synthesizer parameters in prediction order (index 0--15). \textit{Normalisation} is the transform applied to raw Vital values before min-max scaling to $[0,1]$: envelope times and filter cutoff span several orders of magnitude and are $\log_{10}$-transformed first. \textit{Discrete} indicates parameters that Vital internally quantizes to a finite grid; sampled values are snapped to the nearest valid point during preset generation. \textit{Loss wt.} is the group weighting applied in the training objective. Sustain and release are excluded as both stages are unobservable in a fixed 4-second render with no note-off event.}
\label{tab:parameters}
\end{table}

% ─────────────────────────────────────────────────────────
\subsubsection{Step 2: Timbre-Constrained Preset Generation}

A naive uniform sampling of the 16-dimensional parameter space produces predominantly silent, extremely bright, or otherwise pathological sounds of little musical relevance. To ensure that the dataset covers a representative range of practically meaningful sounds, a \textit{timbre-constrained} generation strategy is employed. Five timbre categories are defined, each associated with a set of parameter sub-ranges that constrain the sampling to musically coherent regions of the parameter space. The categories and their MIDI pitch ranges are summarized in Table~\ref{tab:timbre_categories}.

\begin{table}[H]
\centering
\renewcommand{\arraystretch}{1.3}
\begin{tabular}{lcccc}
\hline
\textbf{Timbre} & \textbf{Min MIDI} & \textbf{Max MIDI} & \textbf{Notes} & \textbf{Presets} \\
\hline
bass   & C0 (12) & C3 (48) & Low register, slow envelope      & 6{,}000  \\
lead   & C3 (48) & C6 (84) & Mid-high register, bright timbre  & 6{,}000  \\
pad    & C2 (36) & C5 (72) & Sustained, slow attack            & 6{,}000  \\
percs  & C1 (24) & C6 (84) & Short attack, noise-heavy         & 6{,}000  \\
random & C0 (12) & C6 (84) & Unconstrained full range          & 20{,}000 \\
\hline
\textbf{Total} & & & & \textbf{44{,}000} \\
\hline
\end{tabular}
\vspace{0.3cm}
\caption{Timbre categories used for preset generation, with their MIDI pitch ranges and preset counts.}
\label{tab:timbre_categories}
\end{table}

\noindent For each timbre, a JSON range file defines per-parameter minimum and maximum bounds that enforce the target timbre. For example, the \texttt{percs} category enforces a very short amplitude attack ($\leq 10$ ms), a moderate decay range, and a high noise/sample oscillator level, ensuring that any random draw within those bounds produces a percussive-sounding result. Each preset is stored as a record in a per-timbre JSONL file, with checkpoint-aware deduplication allowing interrupted generation runs to be safely resumed.

% ─────────────────────────────────────────────────────────
\subsubsection{Step 3: Audio Rendering}

Each preset is rendered to a 44.1 kHz mono WAV file via Pedalboard, which loads Vital as a VST3 plugin, sets all 16 parameters programmatically, and triggers a MIDI note-on event without a corresponding note-off. The clip duration is fixed at 4 seconds. The absence of a note-off event means that the release stage of the amplitude envelope is never triggered, which is why release is excluded from the 16 studied parameters. The sustain stage is similarly excluded, as the envelope decay is set long enough relative to the 4-second clip to produce effectively sustained sounds without requiring an explicit sustain parameter.

\par \vspace{0.3cm}
In V1 and V2, each preset is rendered at a single MIDI pitch sampled uniformly from the timbre's valid range. In V3, each preset is rendered at three fixed MIDI pitches chosen to cover the low, mid, and high registers of the timbre's range, expanding the dataset approximately three-fold. A SHA-256 hash of the parameter vector and MIDI pitch is used to track already-rendered presets, enabling safe resumption of interrupted runs. Rendered audio paths and ground-truth parameters are written to a combined JSONL file with one record per clip.

% ─────────────────────────────────────────────────────────
\subsubsection{Step 4: Silence Filtering}

Certain extreme parameter combinations cause Vital to produce no audible output — for example, when both oscillators are muted and the sample level is zero. These silent clips are detected and removed using a hybrid energy-based rule applied to each rendered WAV file.

\par \vspace{0.3cm}
The audio signal $x(n)$ is divided into short frames of length $L$ with hop $H$. The RMS energy of the $i$-th frame is computed as:

\begin{equation}
E_i = \sqrt{\frac{1}{L} \sum_{n=0}^{L-1} x^2(iH + n)}
\label{eq:rms}
\end{equation}

\noindent and converted to decibels:

\begin{equation}
E_i^{\text{dB}} = 20 \cdot \log_{10}(E_i + \epsilon)
\label{eq:rms_db}
\end{equation}

\noindent where $\epsilon$ is a small floor value to avoid $\log(0)$. Two detection rules are then applied:

\begin{equation}
R = \frac{1}{N_f} \sum_{i=1}^{N_f} \mathbf{1}\!\left[E_i^{\text{dB}} > \theta_{\text{rms}}\right]
\label{eq:active_ratio}
\end{equation}

\begin{equation}
P^{\text{dB}} = 20 \cdot \log_{10}\!\left(\max_{n}\,|x(n)|\right)
\label{eq:peak_db}
\end{equation}

\noindent where $R$ is the \textit{active ratio} — the fraction of frames exceeding the RMS threshold $\theta_{\text{rms}}$ — and $P^{\text{dB}}$ is the peak amplitude in decibels. A clip is retained if either condition is satisfied:

\begin{equation}
\text{keep} \;\Leftrightarrow\; \left(R > \rho_{\min}\right) \;\vee\; \left(P^{\text{dB}} > \theta_{\text{peak}}\right)
\label{eq:keep_rule}
\end{equation}

\noindent The logical OR ensures that transient-heavy sounds such as percussions — which concentrate energy in a small number of frames and would fail the active ratio rule alone — are not incorrectly discarded. The per-timbre thresholds are given in Table~\ref{tab:silence_thresholds}.

\begin{table}[H]
\centering
\renewcommand{\arraystretch}{1.3}
\begin{tabular}{lccc}
\hline
\textbf{Timbre} & $\boldsymbol{\theta_{\text{rms}}}$ \textbf{(dB)} & $\boldsymbol{\rho_{\min}}$ & $\boldsymbol{\theta_{\text{peak}}}$ \textbf{(dB)} \\
\hline
bass   & $-50.0$ & $0.020$ & $-28.0$ \\
lead   & $-50.0$ & $0.020$ & $-28.0$ \\
pad    & $-55.0$ & $0.015$ & $-30.0$ \\
percs  & $-50.0$ & $0.002$ & $-30.0$ \\
random & $-50.0$ & $0.010$ & $-30.0$ \\
\hline
\end{tabular}
\vspace{0.3cm}
\caption{Per-timbre silence detection thresholds used in the hybrid RMS filtering step.}
\label{tab:silence_thresholds}
\end{table}

% ─────────────────────────────────────────────────────────
\subsubsection{Step 5: Mel Spectrogram Extraction}

Each surviving audio clip is converted to a mel spectrogram representation. Given the discrete-time audio signal $x(n)$ at sample rate $f_s = 44{,}100$ Hz, the Short-Time Fourier Transform (STFT) is first computed as:

\begin{equation}
X(t, f) = \sum_{n} x(n)\, w(n - tH)\, e^{-j2\pi fn/N}
\label{eq:stft}
\end{equation}

\noindent where $w$ is a Hann window of length $N$ (the FFT size), $H$ is the hop length, and $t$ indexes time frames. The power spectrogram is then obtained as:

\begin{equation}
P(t, f) = |X(t, f)|^2
\label{eq:power_spec}
\end{equation}

\noindent A bank of $M = 128$ triangular mel filters $\{H_m(f)\}_{m=1}^{M}$, spaced according to the mel scale between $f_{\min} = 20$ Hz and $f_{\max} = 12{,}000$ Hz, is applied to yield the mel spectrogram:

\begin{equation}
S(t, m) = 10 \cdot \log_{10}\!\left(\sum_{f} H_m(f)\, P(t, f) + \epsilon\right)
\label{eq:mel}
\end{equation}

\noindent This produces a $128 \times T$ matrix, where $T$ is the number of time frames.

\par \vspace{0.3cm}
\noindent \textbf{V1 — Single-Scale.} A single mel spectrogram is extracted using the mid-scale FFT configuration (Table~\ref{tab:mel_scales}) and stored as a \texttt{.npy} file.

\par \vspace{0.3cm}
\noindent \textbf{V2 — Multi-Scale.} Three mel spectrograms are extracted per clip at fine, mid, and coarse resolutions. As the FFT window size increases, spectral resolution improves at the cost of temporal resolution, and vice versa. Table~\ref{tab:mel_scales} lists the parameters of each scale and the type of information each captures.

\begin{table}[H]
\centering
\renewcommand{\arraystretch}{1.3}
\begin{tabular}{lcccp{4.5cm}}
\hline
\textbf{Scale} & \textbf{n\_fft} & \textbf{Hop} & \textbf{Raw Frames ($T$)} & \textbf{Captures} \\
\hline
Fine   & 512   & 128   & $\approx$1{,}378 & Attack transients, onset details \\
Mid    & 2{,}048 & 512   & $\approx$345     & Timbral balance, envelope shape  \\
Coarse & 8{,}192 & 2{,}048 & $\approx$87      & Harmonic structure, pitch content \\
\hline
\multicolumn{5}{l}{\small All scales: 128 mel bins, $f_{\min}=20$ Hz, $f_{\max}=12{,}000$ Hz, $f_s=44{,}100$ Hz.}\\
\multicolumn{5}{l}{\small All inputs padded or truncated to 224 time frames for ConvNeXt/AST.}\\
\hline
\end{tabular}
\vspace{0.3cm}
\caption{Multi-scale mel spectrogram configurations used in V2 and V3.}
\label{tab:mel_scales}
\end{table}

\noindent The three spectrograms are stacked along the channel dimension to form a three-channel input tensor of shape $(3, 128, 224)$, as illustrated in Figure~\ref{fig:multiscale_stack}. This is directly compatible with the three-channel (RGB) first-layer weight initialization of ConvNeXt and ViT backbones.

\begin{figure}[H]
\centering
\begin{tikzpicture}[
  spec/.style={draw=gray!60, fill=blue!6, minimum width=3.8cm, minimum height=1.1cm,
               align=center, font=\small},
  label/.style={font=\small, align=center},
  arrow/.style={-{Stealth[length=4mm]}, thick, gray!50}
]

% Three spectrogram slabs
\node[spec, fill=orange!12] (fine)   at (0, 2.6) {Fine\\$n_{\text{fft}}=512$};
\node[spec, fill=blue!10]   (mid)    at (0, 1.3) {Mid\\$n_{\text{fft}}=2048$};
\node[spec, fill=teal!10]   (coarse) at (0, 0.0) {Coarse\\$n_{\text{fft}}=8192$};

\node[label] at (0, -0.7) {$128\times224$ each};

% Brace / arrow
\draw[arrow] (2.2, 1.3) -- (3.4, 1.3);

% Stacked tensor box
\draw[fill=gray!8, draw=gray!60, rounded corners=4pt]
  (3.5, 0.3) rectangle (5.5, 2.3);
\node[font=\small, align=center] at (4.5, 1.3)
  {$(3, 128, 224)$\\\small tensor};

% Channel labels inside
\draw[fill=orange!12, draw=gray!40] (3.6, 1.75) rectangle (5.4, 2.15);
\draw[fill=blue!10,   draw=gray!40] (3.6, 1.25) rectangle (5.4, 1.65);
\draw[fill=teal!10,   draw=gray!40] (3.6, 0.45) rectangle (5.4, 0.85);
\node[font=\tiny] at (4.5, 1.95) {ch 0: Fine};
\node[font=\tiny] at (4.5, 1.45) {ch 1: Mid};
\node[font=\tiny] at (4.5, 0.65) {ch 2: Coarse};

\end{tikzpicture}
\vspace{0.3cm}
\caption{Three mel spectrograms at fine, mid, and coarse resolutions stacked into a single three-channel input tensor of shape $(3, 128, 224)$.}
\label{fig:multiscale_stack}
\end{figure}

% ─────────────────────────────────────────────────────────
\subsubsection{Step 6: Dataset Splitting and Normalization}

\subsubsubsection{Stratified Split}
To ensure that each data split contains a representative distribution of both timbre categories and pitch registers, a stratified splitting strategy is employed. Each sample is assigned a bucket key:

\begin{equation}
\text{bucket} = \left(\text{timbre},\; \left\lfloor \frac{\text{midi\_note}}{12} \right\rfloor\right)
\label{eq:bucket}
\end{equation}

\noindent Within each bucket, samples are deterministically shuffled using a seeded random number generator and then divided in an 80 / 10 / 10 ratio into training, validation, and test sets. Five independent splits are generated using different random seeds, producing five-fold cross-validation splits for robust evaluation. In V3, the split is performed at the \textit{preset level}: all three pitch renderings of the same preset are assigned to the same split, preventing any form of data leakage between training and evaluation sets.

\subsubsubsection{Normalization}
Three normalization schemes are applied, each fitted exclusively on training data and applied identically to validation and test sets to prevent data leakage.

\par \vspace{0.3cm}
\textit{Mel spectrograms} are Z-score normalized per scale:

\begin{equation}
\hat{S}(t, m) = \frac{S(t, m) - \mu_{\text{train}}}{\sigma_{\text{train}}}
\label{eq:zscore}
\end{equation}

\noindent \textit{Synthesizer parameters} are min-max normalized per parameter to $[0, 1]$:

\begin{equation}
\hat{x}_k = \frac{x_k - x_k^{\min}}{x_k^{\max} - x_k^{\min}}, \quad k = 1, \ldots, 16
\label{eq:minmax}
\end{equation}

\noindent \textit{MIDI pitch} is normalized to $[0, 1]$ using the fixed global range $[12, 84]$:

\begin{equation}
p = \frac{\text{midi\_note} - 12}{84 - 12}
\label{eq:pitch_norm}
\end{equation}

\noindent This normalized scalar $p \in [0, 1]$ is then passed through a sinusoidal embedding to produce a rich 128-dimensional Fourier feature vector:

\begin{equation}
\boldsymbol{\phi}(p) = \left[\sin(\omega_0 p),\, \cos(\omega_0 p),\, \sin(\omega_1 p),\, \cos(\omega_1 p),\, \ldots,\, \sin(\omega_{63} p),\, \cos(\omega_{63} p)\right]
\label{eq:sinusoidal}
\end{equation}

\noindent where the frequencies $\omega_k = 2^k \cdot \pi$ are log-spaced from $\omega_0 = \pi$ to $\omega_{63} = 2^{63}\pi$. This avoids float32 overflow from naive power scaling and provides the model with a dense, multi-frequency representation of the pitch scalar. The final dataset sizes after filtering and splitting are summarized in Table~\ref{tab:dataset_sizes}.

\begin{table}[H]
\centering
\renewcommand{\arraystretch}{1.3}
\begin{tabular}{lcccccc}
\hline
\textbf{Version} & \textbf{Presets} & \textbf{Pitches / Preset} & \textbf{Total Samples} & \textbf{Train} & \textbf{Val} & \textbf{Test} \\
\hline
V1 / V2 & $\approx$42{,}520 & 1 & $\approx$42{,}520 & $\approx$34{,}016 & $\approx$4{,}252 & $\approx$4{,}252 \\
V3      & $\approx$42{,}520 & 3 & $\approx$127{,}000 & $\approx$101{,}600 & $\approx$12{,}700 & $\approx$12{,}700 \\
\hline
\end{tabular}
\vspace{0.3cm}
\caption{Dataset sizes across experimental versions. V3 splits are performed at preset level to prevent leakage.}
\label{tab:dataset_sizes}
\end{table}

% ─────────────────────────────────────────────────────────
\subsection{Technologies}

Table~\ref{tab:technologies} lists the core tools and libraries used across the data generation, model training, and evaluation stages of this project.

\begin{table}[H]
\centering
\renewcommand{\arraystretch}{1.3}
\begin{tabular}{p{3.2cm} p{10cm}}
\hline
\textbf{Tool / Library} & \textbf{Role} \\
\hline
Vital                   & Open-source wavetable VST3 synthesizer; audio generation engine and source of ground-truth parameters \\
Pedalboard (Spotify)    & Python VST plugin host; used to load Vital, set parameters programmatically, and render audio clips \\
Librosa                 & Audio analysis: mel spectrogram extraction, RMS energy computation, spectral descriptor calculation \\
SoundFile               & Low-level WAV file reading and writing \\
NumPy                   & Numerical operations, array manipulation, and preset JSONL handling \\
PyTorch                 & Model definition, training loop, loss computation, and inference \\
timm                    & Pre-trained vision model zoo; provides ConvNeXt and ViT/AST backbones with ImageNet-22k weights \\
Python 3.10             & Full pipeline orchestration across all six data generation steps and model training \\
CUDA / AMP              & GPU-accelerated training on NVIDIA RTX 4090 with automatic mixed precision \\
\hline
\end{tabular}
\vspace{0.3cm}
\caption{Tools and libraries used in the data generation, training, and evaluation pipeline.}
\label{tab:technologies}
\end{table}

